\chapter{Probability and Predictive Reconstruction}
\label{s:probability}

\section{The conditional relative-entropy identity}

The disintegrations used here are regular conditional laws on standard Borel spaces \citep{Kallenberg2021}.

Let $\mu,\nu$ be probability laws on a product $X\times Y$ of standard Borel spaces. Write their marginals as $\mu_Y,\nu_Y$ and their conditional laws as $\mu^y,\nu^y$. Then
\begin{equation}
\KL(\mu\Vert\nu)
=\KL(\mu_Y\Vert\nu_Y)
+\int\KL(\mu^y\Vert\nu^y)\mu_Y(dy).
\label{eq:s-chain-rule}
\end{equation}
The identity holds in $[0,\infty]$, both sides being infinite together. Suppose first $\mu\ll\nu$ and put $f=d\mu/d\nu$, a jointly measurable function. Then $a(y)=\int f(x,y)\nu^y(dx)$ is a version of $d\mu_Y/d\nu_Y$, and $b=f/a$ on $\{a>0\}$ is jointly measurable with $\mu^y=b(\cdot,y)\nu^y$ for $\mu_Y$-almost every $y$. Starting from $f$ in this way avoids choosing a jointly measurable version of the conditional densities separately. The logarithm splits into $\log a+\log b$. The negative parts of $x\log x$ are bounded by $1/e$, so the negative parts of the two logarithms have finite $\mu$-integrals and the split integrates in $[0,\infty]$, giving \cref{eq:s-chain-rule}. If $\mu$ is not absolutely continuous with respect to $\nu$, then either $\mu_Y\not\ll\nu_Y$ or $\mu^y\not\ll\nu^y$ on a set of positive $\mu_Y$-measure, and both sides are infinite. This argument underlies the predictive and reconstruction identities in the main text.

\section{Different predictive mechanisms}

Let $s,t$ be laws on a model space and let $K,L$ be normalized kernels into a common record space. Suppose the following joint divergence is finite. Define
\begin{equation}
\mu(dm,do)=s(dm)K(m,do),
\qquad \nu(dm,do)=t(dm)L(m,do).
\end{equation}
Disintegration first by $m$ and then by $o$ gives the exact equality
\begin{equation}
\begin{aligned}
\KL(s\Vert t)+\int\KL(K(m,\cdot)\Vert L(m,\cdot))s(dm)
={}&\KL(sK\Vert tL)\\
&+\int\KL(\mu^o\Vert\nu^o)(sK)(do).
\end{aligned}
\label{eq:s-different-evaluators}
\end{equation}
This is two uses of \cref{eq:s-chain-rule} on one joint law. It displays explicitly the distinction contributed by the evaluators. Equal model laws do not remove that term.

For a bijective change of model presentation $\Theta$, replace $t$ by $\Theta_\#t$ and replace $L$ by $L(\Theta^{-1}m,\cdot)$. Both sides then compare laws in one presentation. Measurable record relabeling must be applied consistently as well.

\section{A conditional reconstruction on four points}

Let $E=\{0,1\}^2$, write its coordinates as $(H,V)$, and transmit only $H$. Choose a sender law with
\begin{equation}
q(H=1)=\frac12,
\qquad q(V=H\mid H)=1.
\end{equation}
Let the receiver's prior $p$ make $H,V$ independent uniform bits. The communicated marginal is uniform, so reconstruction using $p(V\mid H)$ returns the uniform law on all four points. Its transmitted marginal is exactly correct, but it discards the sender's perfect dependence. Direct evaluation gives
\begin{equation}
\KL(q\Vert\widehat q)
=2\left(\frac12\log\frac{1/2}{1/4}\right)=\log 2.
\end{equation}
The conditional divergence at each value of $H$ is also $\log2$, verifying the conditional identity exactly. A receiver whose prior has the same conditional relation $V=H$ instead reconstructs $q$ from the transmitted marginal. The difference comes from its prior conditional, not from a larger message. The reconstruction is Jeffrey's rule of conditioning the receiver's prior on the transmitted marginal \citep{DiaconisZabell1982}, and it is the law closest to that prior, in relative entropy, among all laws with the transmitted marginal.

\section{A shared source and two routes}

Let $X$ be a uniform bit. Two routes return $Y_1=X$ and $Y_2=X\oplus b$, where $b\in\{0,1\}$ is a fixed route mark and $\oplus$ is addition modulo two. The comparison $O=Y_1\oplus Y_2$ equals $b$ with probability one. If the second route instead uses an independent uniform bit $X'$, the same comparison is uniform and reveals no route mark. The marginal law arriving on each route is uniform in both experiments. Thus the shared-source dependence, not those marginals, determines the comparison outcome. A record model intended to inspect a loop must specify which occurrences refer to the same source. In general, a functional of the pair of route laws cannot distinguish two loop maps $H,H'$ with the same pushforward of the source law, whereas the joint law of a shared sample and its image along the loop determines the loop map almost everywhere; the route mark is the simplest instance.

\section{Predictive sufficiency for a selected response}

Let $X=(A,B)$ be two independent uniform bits, let the encoder retain $U=A$, and take no side information. For target $Y=A$, one has $I(Y;X\mid U)=0$. For target $Y=B$, the same representation gives $I(Y;X\mid U)=\log2$. Perfect predictive retention is therefore a statement about a specified target. If a later interaction reads $B$, the first sufficiency statement does not supply its response. This elementary distinction persists when $X$ is a complete configuration of law-valued agents and $Y$ is a boundary process.
